\section{Discussion and Empirical Implications}
\label{sec:discussion}

The model is stylized, but it generates several empirical predictions that distinguish the incidence channel from a conventional procurement-efficiency effect.

First, under the baseline incidence and DRHR conditions, greater supplier competition should weaken a jurisdiction's willingness to accept a supplier-side policy role when locally captured supplier rents are economically or politically salient. The relevant outcome is not whether procurement becomes cheaper conditional on a project. The prediction concerns whether the complementary activity is selected at all. Empirically, one would therefore look for changes in regional specialization, local-content support, supplier-development programs, or willingness to host complementary production stages as supplier concentration changes.

Second, the effect should be heterogeneous in local incidence. Regions whose governments internalize a larger share or weight of supplier-side surplus should have a larger effective $\rho$. They should therefore display a stronger negative relationship between supplier competition and willingness to accept the $D$ role. Conversely, when less supplier-side surplus enters the regional government's objective, $\rho$ is smaller and the rent-incidence mechanism is weaker in levels. If the jurisdiction also captures buyer-side gains---represented by $\sigma>0$ in Section \ref{sec:robustness}---the competition effect on local incentives can weaken or reverse.

Third, demand for intergovernmental coordination can increase precisely when procurement performance improves. In the model, a larger $N$ reduces expected production cost and improves total real surplus conditional on differentiated production. At the same time, it expands the parameter region in which decentralized governments duplicate $U$. This suggests an empirical distinction between \emph{within-project efficiency} and \emph{across-region policy coordination}. Better procurement outcomes need not imply better decentralized industrial-policy composition.

These predictions relate to, but differ from, the literature on procurement preferences. Bid preferences and set-asides deliberately alter the rules faced by favored firms \citep{Marion2007,AtheyCoeyLevin2013}. Our mechanism operates even when the procurement mechanism remains neutral: the distribution of supplier rents changes because the supplier market is more competitive. Nor is the mechanism one of local governments bidding against each other for firms, as in \citet{Slattery2025}. Here governments choose policy roles, and competition occurs among firms in the continuation market.

The theory also clarifies the appropriate policy interpretation. The model does not imply that governments should restrict supplier competition to preserve local rents. Doing so would sacrifice procurement efficiency and may introduce familiar favoritism and entry distortions. The result instead identifies a coordination problem: if competition erodes the local return to a socially valuable complementary role, the natural response may involve intergovernmental coordination, transfers, or institutional arrangements that reallocate some of the project surplus. Those instruments are deliberately not modeled here. Introducing them would change the policy mechanism and is left for future work.

Finally, the binary role choice should not be read as a literal claim that real regional policy is all-or-nothing. Its purpose is to isolate the equilibrium-regime effect. A continuous portfolio version would require specifying how procurement surplus scales with partial supplier-side policy and would move the paper toward the companion continuous-composition model. The present binary structure makes the market-structure comparative static and the disappearance of differentiated Nash equilibria exact.
